\documentclass[11pt,reqno]{amsart}

\usepackage[T1]{fontenc}
\usepackage{lmodern}
\usepackage{microtype}
\usepackage{mathtools,amssymb,amsfonts}
\usepackage{booktabs}
\usepackage{enumitem}
\usepackage[margin=1.08in]{geometry}
\usepackage{xcolor}
\definecolor{linkblue}{HTML}{244E73}
\definecolor{citegreen}{HTML}{315B45}
\usepackage[
  colorlinks=true,
  linkcolor=linkblue,
  citecolor=citegreen,
  urlcolor=linkblue,
  pdfauthor={Working draft},
  pdftitle={Makeev's regular-simplex conjecture: positive results and a counterexample in dimension four}
]{hyperref}
\usepackage[nameinlink,noabbrev]{cleveref}

\allowdisplaybreaks
\setlist[enumerate]{leftmargin=2.2em,itemsep=0.25em,topsep=0.4em}
\setlist[itemize]{leftmargin=2.0em,itemsep=0.25em,topsep=0.4em}

\newtheorem{theorem}{Theorem}[section]
\newtheorem{proposition}[theorem]{Proposition}
\newtheorem{lemma}[theorem]{Lemma}
\newtheorem{corollary}[theorem]{Corollary}
\theoremstyle{definition}
\newtheorem{definition}[theorem]{Definition}
\theoremstyle{remark}
\newtheorem{remark}[theorem]{Remark}

\newcommand{\R}{\mathbb R}
\newcommand{\C}{\mathbb C}
\newcommand{\one}{\mathbf 1}
\newcommand{\SO}{\mathrm{SO}}
\newcommand{\OO}{\mathrm{O}}
\newcommand{\proj}{\operatorname{proj}}
\newcommand{\Sym}{\operatorname{Sym}}
\newcommand{\dd}{\mathop{}\!\mathrm d}
\newcommand{\eps}{\varepsilon}
\newcommand{\cE}{\mathcal E}
\newcommand{\cC}{\mathcal C}
\newcommand{\cT}{\mathcal T}
\newcommand{\cH}{\mathcal H}
\newcommand{\ip}[2]{\left\langle #1,#2\right\rangle}
\newcommand{\norm}[1]{\left\lVert #1\right\rVert}
\newcommand{\abs}[1]{\left\lvert #1\right\rvert}

\title[Makeev's conjecture in dimension four]{Makeev's regular-simplex conjecture:\\positive results and a counterexample in dimension four}
\author{Working draft}
\thanks{This manuscript reorganizes and simplifies the supplied arguments. It makes no claim of priority.}
\date{August 2026}

\begin{document}

\begin{abstract}
Let
\[
V=\one^\perp\subset \R^{d+1},\qquad
u_{ij}=\frac{e_i-e_j}{\sqrt2}\in S(V),\qquad 0\le i<j\le d,
\]
be the normalized roots of type $A_d$.  We study whether every smooth odd
$f\colon S(V)\to\R$ agrees with a linear functional on a rotated copy of these
roots.  A single variational identity identifies the cycle defect of a cubic
with the gradient of its correlation with $\sum x_i^3$.  It yields the desired
conclusion for all spherical harmonics of degrees $1$ and $3$, for a nonempty
$C^0$-open set of arbitrary odd functions, and for all zonal odd functions.

In dimension $4$ we then give an explicit counterexample.  Identifying
$\R^4\cong\C^2$, for every sufficiently small nonzero real $\eps$ the function
\[
\operatorname{Re}(z_1^2\overline z_2)
+\eps\left(\abs{z_1}^2-\frac35\right)
 \operatorname{Im}(z_1^2\overline z_2),\qquad (z_1,z_2)\in S^3,
\]
does not agree with any linear functional on any rotated $A_4$ root system.
The two summands are spherical harmonics of degrees $3$ and $5$.  The main
proof is reduced to a symmetry identity, a compact obstruction lemma, and one
explicit algebraic transversality proposition; the complete algebraic
verification is placed in an appendix.
\end{abstract}

\maketitle

\section{Introduction}

Makeev conjectured that the dual difference body of a regular simplex is a
universal cover for sets of diameter one \cite{Makeev1994}.  Kuperberg showed
that the constant-width formulation is equivalent to the following
function-theoretic assertion: every continuous odd function on $S^{d-1}$
agrees with a linear function on a rotated copy of the $A_d$ root system.  He
proved the assertion for $d=3$ and showed that the homological obstruction used
there vanishes in every higher dimension \cite{Kuperberg1999}.  A related
equivariant proof in dimension three was given by Hausel, Makai, and
Sz\H{u}cs \cite{HauselMakaiSzucs2000}.  Chen's July 2026 preprint still records
the regular-$A_d$ universal-cover statement as Makeev's conjecture
\cite{Chen2026}.

We normalize the roots throughout:
\begin{equation}\label{eq:normalized-roots}
V=\one^\perp\subset\R^m,\qquad m=d+1,\qquad
u_{ij}=\frac{e_i-e_j}{\sqrt2}\in S(V).
\end{equation}
This normalization is essential because the function is defined on the unit
sphere.  The root configuration has an orientation-reversing symmetry, for
example the restriction to $V$ of a transposition of two coordinate vectors.
Consequently, replacing $\OO(V)$ by $\SO(V)$ does not change the problem.

The first part of the paper records three positive results.  The main device is
a gradient identity for cubics, closely related in spirit to variational
orthogonal tensor diagonalization \cite{MartinVanLoan2008}.  The second part
gives a negative answer in dimension four.  Its perturbative structure is
similar to the compactness-plus-infinitesimal-obstruction strategy that appears
in Karasev's counterexample to another conjecture of Makeev
\cite{Karasev2016}, although the obstruction here is different and completely
explicit.

\begin{theorem}[Main counterexample]\label{thm:main-counterexample}
There is $\eps_0>0$ such that, for every real $\eps$ with
$0<\abs\eps<\eps_0$, the smooth odd function
\begin{equation}\label{eq:main-counterexample}
\widetilde f_\eps(z_1,z_2)
=
\operatorname{Re}(z_1^2\overline z_2)
+\eps\left(\abs{z_1}^2-\frac35\right)
 \operatorname{Im}(z_1^2\overline z_2)
\end{equation}
on $S^3\subset\C^2$ has no $A\in\SO(4)$ and $b\in\R^4$ satisfying
\[
\widetilde f_\eps(A\nu_{ij})=b\cdot A\nu_{ij}
\qquad(0\le i<j\le4).
\]
\end{theorem}

\section{Edge flows and the test map}\label{sec:test-map}

Let $\cE$ be the space of skew edge labels on the complete graph $K_m$:
\[
\cE=\{Y\in\R^{m\times m}:Y^T=-Y\},
\qquad
\ip{X}{Y}_{\cE}=\sum_{i<j}X_{ij}Y_{ij}.
\]
For $c\in V$, define
\[
(\delta c)_{ij}=\frac{c_i-c_j}{\sqrt2}.
\]
These are exactly the restrictions of linear functions on $V$ to the roots
\eqref{eq:normalized-roots}.  Put
\begin{equation}\label{eq:cycle-space}
\cC=\{X\in\cE:X\one=0\}.
\end{equation}
This is the divergence-free or cycle space.  The orthogonal splitting
$\cE=\delta V\oplus\cC$ is the complete-graph instance of combinatorial Hodge
theory for edge flows \cite{JiangLimYaoYe2011}.

\begin{lemma}[Exact edge labels]\label{lem:exact-labels}
For $Y\in\cE$, the following are equivalent.
\begin{enumerate}
\item There is $c\in V$ such that
\[
Y_{ij}=\frac{c_i-c_j}{\sqrt2}.
\]
\item Every triangle sum vanishes:
\[
Y_{ij}+Y_{jk}+Y_{ki}=0.
\]
\item $Y$ is orthogonal to $\cC$.
\item With $\Pi=I-m^{-1}\one\one^T$, one has
\[
\Pi Y\Pi=0.
\]
\end{enumerate}
\end{lemma}

\begin{proof}
The first condition implies the second by telescoping.  Conversely, if all
triangle sums vanish, set $c_0=0$ and $c_i=\sqrt2Y_{i0}$.  The triangle
$(i,j,0)$ gives $Y_{ij}=(c_i-c_j)/\sqrt2$; subtracting the mean of the $c_i$
places $c$ in $V$.

If $X\one=0$, then
\[
\sum_{i<j}X_{ij}(c_i-c_j)=\sum_i c_i\sum_jX_{ij}=0,
\]
so exact labels are orthogonal to $\cC$.  The divergence map
$X\mapsto X\one$ maps $\cE$ onto $V$: for $v\in V$, the skew matrix
$X_{ij}=v_i-v_j$ satisfies $X\one=mv$.  Hence
\[
\dim\cC=\binom m2-(m-1),
\]
and the orthogonal complement of $\cC$ has the same dimension as $\delta V$.
This proves the equivalence of the first three conditions.

For the matrix criterion, an exact label has the form
$Y=\rho\one^T-\one\rho^T$, so $\Pi Y\Pi=0$.  Conversely, suppose
$\Pi Y\Pi=0$, and write $J_0=m^{-1}\one\one^T$.  Since $Y$ is skew,
$J_0YJ_0=0$, and hence
\[
Y=J_0Y+YJ_0.
\]
If $\rho_i=m^{-1}\sum_jY_{ij}$, then
$(YJ_0)_{ij}=\rho_i$ and $(J_0Y)_{ij}=-\rho_j$.  Thus
$Y_{ij}=\rho_i-\rho_j$, which is exact after taking
$c_i=\sqrt2\rho_i$ and subtracting the mean.
\end{proof}

The space $\cC$ is naturally $\mathfrak{so}(V)$: extend a skew operator on
$V$ by zero on $\R\one$.  For an odd continuous function
$f\colon S(V)\to\R$ and $A\in\SO(V)$, define
\begin{equation}\label{eq:test-map}
Y_f(A)_{ij}=f(A\nu_{ij}),\qquad
\cT_f(A)=\proj_{\cC}Y_f(A).
\end{equation}

\begin{corollary}[Test-map formulation]\label{cor:test-map}
There are $A\in\SO(V)$ and $b\in V$ such that
\[
f(A\nu_{ij})=b\cdot A\nu_{ij}\qquad(i<j)
\]
if and only if $\cT_f(A)=0$ for some $A\in\SO(V)$.
\end{corollary}

\begin{proof}
If $\cT_f(A)=0$, then \cref{lem:exact-labels} gives $c\in V$ with
$f(A\nu_{ij})=c\cdot\nu_{ij}$.  Taking $b=Ac$ gives the stated identity.
The converse is immediate.
\end{proof}

Notice the dimension equality
\[
\dim\SO(d)=\binom d2=\dim\cC.
\]
It is the source of the degree and Euler-class approaches in the earlier
literature \cite{Kuperberg1999,HauselMakaiSzucs2000}.

\section{The cubic gradient identity}\label{sec:cubic}

Define the distinguished cubic
\begin{equation}\label{eq:distinguished-cubic}
p(x)=\sum_{i=0}^{d}x_i^3,
\qquad x=(x_0,\dots,x_d)\in V.
\end{equation}
If $Q$ is a cubic on $V$, let $T_Q$ be its symmetric trilinear polarization,
so that $Q(x)=T_Q(x,x,x)$.  Extend $T_Q$ to $\R^m$ by projecting each argument
to $V$, and use the invariant tensor inner product
\[
\ip{Q}{R}=\sum_{a,b,c=0}^{d}(T_Q)_{abc}(T_R)_{abc}.
\]
The group acts on cubics by $(g\cdot Q)(x)=Q(g^{-1}x)$.

\begin{lemma}[Cubic edge identity]\label{lem:cubic-edge-identity}
For every cubic $Q$ on $V$ and every $X\in\cC=\mathfrak{so}(V)$,
\begin{equation}\label{eq:cubic-edge-identity}
\ip{Q}{X\cdot p}
=2\sqrt2\sum_{i<j}X_{ij}Q(\nu_{ij}).
\end{equation}
\end{lemma}

\begin{proof}
The infinitesimal action gives
\[
(X\cdot p)(x)
=\left.\frac{\dd}{\dd t}\right|_{t=0}p(e^{-tX}x)
=-3\sum_{i,j}X_{ij}x_i^2x_j,
\]
so
\begin{equation}\label{eq:tensor-side}
\ip{Q}{X\cdot p}=-3\sum_{i,j}X_{ij}(T_Q)_{iij}.
\end{equation}
On the other hand,
\[
Q(e_i-e_j)
=(T_Q)_{iii}-3(T_Q)_{iij}+3(T_Q)_{ijj}-(T_Q)_{jjj}.
\]
After multiplying by $X_{ij}$ and summing, the pure diagonal terms vanish
because $X\one=0$.  Interchanging $i$ and $j$ in the remaining second term
shows
\[
\sum_{i<j}X_{ij}Q(e_i-e_j)
=-3\sum_{i,j}X_{ij}(T_Q)_{iij}.
\]
Comparison with \eqref{eq:tensor-side}, followed by
$Q(e_i-e_j)=2\sqrt2Q(\nu_{ij})$, proves the identity.
\end{proof}

For a cubic $Q$, set
\[
\Psi_Q(A)=\ip{Q}{A\cdot p},\qquad A\in\SO(V).
\]
Identify $T_A\SO(V)$ with $\{AX:X\in\mathfrak{so}(V)\}$ using the curve
$Ae^{tX}$.

\begin{proposition}[Gradient form of the test map]\label{prop:gradient}
For every $A\in\SO(V)$ and $X\in\cC$,
\begin{equation}\label{eq:gradient-identity}
\dd\Psi_Q(A)[AX]
=2\sqrt2\,\ip{\cT_Q(A)}{X}_{\cE}.
\end{equation}
\end{proposition}

\begin{proof}
By invariance of the tensor inner product and
\cref{lem:cubic-edge-identity},
\begin{align*}
\dd\Psi_Q(A)[AX]
&=\ip{Q}{A\cdot(X\cdot p)}
 =\ip{A^{-1}\cdot Q}{X\cdot p}\\
&=2\sqrt2\sum_{i<j}X_{ij}Q(A\nu_{ij})
 =2\sqrt2\,\ip{\cT_Q(A)}{X}_{\cE}.
\end{align*}
The last equality uses $X\in\cC$, so projection onto $\cC$ does not change
the pairing.
\end{proof}

\begin{theorem}[Cubics plus linear functions]\label{thm:cubics}
Let $d\ge2$.  If
\[
f(x)=Q(x)+\ell(x),\qquad x\in S(V),
\]
where $Q$ is homogeneous cubic and $\ell$ is linear, then there are
$A\in\SO(V)$ and $b\in V$ such that
\[
f(A\nu_{ij})=b\cdot A\nu_{ij}\qquad(i<j).
\]
Equivalently, the conclusion holds for every odd function whose spherical
harmonic expansion has only degrees $1$ and $3$.
\end{theorem}

\begin{proof}
The compact group $\SO(V)$ contains a critical point $A$ of $\Psi_Q$.
By \eqref{eq:gradient-identity}, $\cT_Q(A)=0$.  Thus
\cref{lem:exact-labels} gives $c\in V$ such that
$Q(A\nu_{ij})=c\cdot\nu_{ij}$.  If $\ell(x)=a\cdot x$, then with
$b=Ac+a$,
\[
b\cdot A\nu_{ij}=Q(A\nu_{ij})+\ell(A\nu_{ij}).
\]
\end{proof}

We will also use a coordinate-free form.  Let
$q_0,\dots,q_d\in\R^d$ be the vertices of a centered regular simplex with
\[
\ip{q_i}{q_j}=\delta_{ij}-\frac1{d+1},
\qquad
u_{ij}=\frac{q_i-q_j}{\sqrt2},
\]
and put $p_q(x)=\sum_i\ip{q_i}{x}^3$.  Transporting
\eqref{eq:cubic-edge-identity} by an isometry gives
\begin{equation}\label{eq:coordinate-free-cubic}
\ip{C}{K\cdot p_q}
=2\sqrt2\sum_{i<j}\ip{q_i}{Kq_j}C(\nu_{ij}).
\end{equation}

\begin{corollary}[Annihilation from a continuous symmetry]\label{cor:symmetry-annihilation}
Let $h_\theta=e^{\theta J}\subset\SO(d)$ be a one-parameter subgroup, and let
$C$ be a cubic satisfying $C\circ h_\theta=C$.  For $g\in\SO(d)$ define
\[
X_C(g)_{ij}=\ip{gq_i}{Jgq_j},
\qquad
\Lambda_g(Y)=\sum_{i<j}X_C(g)_{ij}Y_{ij}.
\]
Then
\begin{equation}\label{eq:global-annihilation}
\Lambda_g\bigl((C(g\nu_{ij}))_{ij}\bigr)=0
\qquad\text{for every }g\in\SO(d).
\end{equation}
\end{corollary}

\begin{proof}
The function $\Psi_C(g)=\ip{C}{g\cdot p_q}$ is invariant under left
multiplication by $h_\theta$.  At $g$, the corresponding tangent vector is
$Jg=gK$, where $K=g^{-1}Jg$.  Applying \eqref{eq:coordinate-free-cubic} to the
derivative in this direction yields
\[
0=2\sqrt2\sum_{i<j}\ip{q_i}{Kq_j}C(g\nu_{ij})
 =2\sqrt2\sum_{i<j}\ip{gq_i}{Jgq_j}C(g\nu_{ij}).
\]
\end{proof}

\section{Two further positive results}\label{sec:positive}

\subsection{A \texorpdfstring{$C^0$}{C0}-open set of arbitrary odd functions}

\begin{theorem}[Uniform stability near \texorpdfstring{$p$}{p}]\label{thm:uniform-stability}
For every $d\ge2$ there exists $\eta_d>0$ such that every continuous odd
$f\colon S(V)\to\R$ satisfying
\[
\norm{f-p}_\infty<\eta_d
\]
has a zero of its test map.  The same is true near every nonzero scalar
multiple or rotation of $p$, and adding a linear function does not affect the
conclusion.
\end{theorem}

\begin{proof}
Since $p(\nu_{ij})=0$, one has $\cT_p(I)=0$.  For
$X\in\cC=\mathfrak{so}(V)$ and $u=\nu_{ij}$,
\[
\dd p_u(Xu)=3\sum_k u_k^2(Xu)_k
=\frac32\bigl((Xu)_i+(Xu)_j\bigr).
\]
Now
\[
(Xu)_i=(Xu)_j=-\frac{X_{ij}}{\sqrt2},
\]
and hence
\begin{equation}\label{eq:derivative-test-p}
D\cT_p(I)[X]=-\frac3{\sqrt2}X.
\end{equation}
Thus $I$ is a nondegenerate zero.  By the inverse function theorem, choose a
small exponential-coordinate ball $B\subset\mathfrak{so}(V)$ on which this is
the only zero.  Its local Brouwer degree is the sign of the determinant in
\eqref{eq:derivative-test-p}, and is therefore nonzero.  Standard homotopy
invariance and stability of local degree under uniform perturbations then
apply \cite{DincaMawhin2021}.

For completeness, if $N=\binom{d+1}{2}$, then uniformly in $A$,
\[
\norm{\cT_f(A)-\cT_p(A)}
\le \sqrt N\,\norm{f-p}_\infty.
\]
Taking $\norm{f-p}_\infty$ smaller than the distance of
$\cT_p(e^X)$ from $0$ on $\partial B$, divided by $2\sqrt N$, keeps the entire
straight-line homotopy away from $0$ on $\partial B$.  Nonzero degree therefore
forces a zero of $\cT_f$ in $B$.
\end{proof}

\subsection{Zonal odd functions}

\begin{theorem}[Zonal functions]\label{thm:zonal}
Let $d\ge2$, $a\in S(V)$, and let $h\colon[-1,1]\to\R$ be continuous and odd.
Then every function
\[
f(x)=h(a\cdot x)+\ell(x),
\]
with $\ell$ linear, satisfies the desired conclusion.
\end{theorem}

\begin{proof}
It is enough to take $\ell=0$.  Choose a nontrivial partition
$\{0,\dots,d\}=S\sqcup T$, with $s=\abs S$, $t=\abs T$, and $s+t=m$.
Define $c\in V$ by
\[
c_i=
\begin{cases}
\alpha=\sqrt{\dfrac{t}{sm}},&i\in S,\\[2mm]
\beta=-\sqrt{\dfrac{s}{tm}},&i\in T.
\end{cases}
\]
Then $\sum_i c_i=0$ and $\norm c=1$.  Choose $A\in\SO(V)$ with $A^Ta=c$.
Thus
\[
a\cdot A\nu_{ij}=\frac{c_i-c_j}{\sqrt2}.
\]
Putting $\kappa=(\alpha-\beta)/\sqrt2$, the values are $0$ within either block
and $\pm\kappa$ across the two blocks.  Since $h$ is odd, define
\[
z_i=
\begin{cases}
\sqrt2\,h(\kappa),&i\in S,\\
0,&i\in T.
\end{cases}
\]
After subtracting its mean, $z\in V$, and
\[
h(a\cdot A\nu_{ij})=\frac{z_i-z_j}{\sqrt2}=z\cdot\nu_{ij}.
\]
Taking $b=Az$ proves the claim.
\end{proof}

\section{The dimension-four counterexample}\label{sec:counterexample}

\subsection{A cyclic regular simplex}

Identify $\R^4$ with $\C^2$, equipped with the real inner product
\[
\ip{z}{w}_{\R}=\operatorname{Re}(z_1\overline w_1+z_2\overline w_2).
\]
Let $\zeta=e^{2\pi i/5}$ and define
\begin{equation}\label{eq:cyclic-simplex}
q_j=\sqrt{\frac25}(\zeta^j,\zeta^{2j}),\qquad j=0,\dots,4.
\end{equation}

\begin{lemma}\label{lem:cyclic-simplex}
The vectors $q_0,\dots,q_4$ are the vertices of a centered regular
$4$-simplex:
\[
\sum_{j=0}^4q_j=0,
\qquad
\ip{q_i}{q_j}_{\R}=
\begin{cases}
\frac45,&i=j,\\[1mm]
-\frac15,&i\ne j.
\end{cases}
\]
Consequently
\begin{equation}\label{eq:a4-cyclic-roots}
\nu_{ij}=\frac{q_i-q_j}{\sqrt2}
\end{equation}
are unit vectors forming an $A_4$ root system.
\end{lemma}

\begin{proof}
The sum vanishes because both $\sum_j\zeta^j$ and $\sum_j\zeta^{2j}$ vanish.
Moreover,
\[
\ip{q_i}{q_j}_{\R}
=\frac25\operatorname{Re}\bigl(\zeta^{i-j}+\zeta^{2(i-j)}\bigr).
\]
For $i\ne j$ the real part is $-1/2$, so the inner product is $-1/5$;
for $i=j$ it is $4/5$.  Hence $\norm{q_i-q_j}^2=2$.
The isometry $e_j-\frac15\one\mapsto q_j$ identifies this model with the
standard one.  Composing with an odd permutation of the vertices if necessary
makes the isometry orientation-preserving.
\end{proof}

Define
\begin{equation}\label{eq:QRH}
W(z)=z_1^2\overline z_2,
\qquad
Q(z)=\operatorname{Re}W(z),
\qquad
R(z)=\operatorname{Im}W(z),
\qquad
H(z)=\abs{z_1}^2R(z).
\end{equation}
These are odd real polynomials.  Consider the one-parameter subgroup
\begin{equation}\label{eq:weighted-circle}
h_\theta(z_1,z_2)=(e^{i\theta}z_1,e^{2i\theta}z_2),
\qquad
J(z_1,z_2)=(iz_1,2iz_2).
\end{equation}
Both $Q$ and $R$ are invariant under $h_\theta$.

For $g\in\SO(4)$ put $p_i=gq_i$ and define
\begin{equation}\label{eq:X-Lambda}
X(g)_{ij}=\ip{p_i}{Jp_j}_{\R},
\qquad
\Lambda_g(Y)=\sum_{i<j}X(g)_{ij}Y_{ij}.
\end{equation}
The matrix $X(g)$ is skew, and
\[
\sum_jX(g)_{ij}=\ip{p_i}{J\sum_jp_j}_{\R}=0,
\]
so $X(g)\in\cC$.  Therefore $\Lambda_g$ annihilates every exact label.
By \cref{cor:symmetry-annihilation},
\begin{equation}\label{eq:QR-annihilation}
\Lambda_g\bigl(Y_Q(g)\bigr)=0,
\qquad
\Lambda_g\bigl(Y_R(g)\bigr)=0
\qquad\text{for every }g\in\SO(4).
\end{equation}

The only special calculation needed in the main argument is the following.
Its proof is given in \cref{app:transversality}.

\begin{proposition}[Algebraic transversality]\label{prop:algebraic-transversality}
If the labels $\bigl(Q(g\nu_{ij})\bigr)_{i<j}$ are exact, then
\begin{equation}\label{eq:transversality}
\Lambda_g\bigl(Y_H(g)\bigr)\ne0.
\end{equation}
\end{proposition}

\subsection{A compact obstruction lemma}

\begin{lemma}[Compact obstruction]\label{lem:compact-obstruction}
Let $M$ be compact, let $E$ be a finite-dimensional normed space, and let
$F_0,F_1\colon M\to E$ be continuous.  Suppose that $x\mapsto\lambda_x\in E^*$
is continuous and
\begin{align}
\lambda_x(F_0(x))&=0 &&\text{for every }x\in M,\label{eq:obstruction-a}\\
F_0(x)=0&\Longrightarrow \lambda_x(F_1(x))\ne0.\label{eq:obstruction-b}
\end{align}
Then there is $\eps_0>0$ such that
\[
F_0(x)+\eps F_1(x)\ne0
\]
for every $x\in M$ and every $0<\abs\eps<\eps_0$.
\end{lemma}

\begin{proof}
Let $Z=F_0^{-1}(0)$.  By compactness and
\eqref{eq:obstruction-b}, there is a neighborhood $U$ of $Z$ on which
$\lambda_x(F_1(x))\ne0$.  On the compact set $M\setminus U$,
\[
a:=\min_{M\setminus U}\norm{F_0}>0.
\]
Let $B=\max_M\norm{F_1}$.  For $0<\abs\eps<a/(B+1)$ the sum cannot vanish on
$M\setminus U$.  If it vanished at $x\in U$, applying $\lambda_x$ and using
\eqref{eq:obstruction-a} would give
$0=\eps\lambda_x(F_1(x))$, a contradiction.
\end{proof}

\subsection{Completion of the proof}

Apply \cref{lem:compact-obstruction} with
\[
M=\SO(4),\qquad E=\cC,
\qquad F_0=\cT_Q,\qquad F_1=\cT_H,
\qquad \lambda_g(Y)=\ip{X(g)}{Y}_{\cE}.
\]
Equation \eqref{eq:QR-annihilation} gives
$\lambda_g(F_0(g))=0$ for all $g$.  If $F_0(g)=0$, then the $Q$-labels are
exact, and \cref{prop:algebraic-transversality} gives
$\lambda_g(F_1(g))\ne0$.  Hence, for all sufficiently small nonzero $\eps$,
\[
\cT_{Q+\eps H}(g)\ne0\qquad\text{for every }g\in\SO(4).
\]
By \cref{cor:test-map}, $Q+\eps H$ is a counterexample.

It remains to obtain the harmonic form stated in
\cref{thm:main-counterexample}.  The cubic $R$ is harmonic, and in $\R^4$ one
checks
\[
\Delta(\abs{z_1}^2R)=12R,
\qquad
\Delta\bigl((\abs{z_1}^2+\abs{z_2}^2)R\bigr)=20R.
\]
Thus
\begin{equation}\label{eq:harmonic-quintic}
H_5(z)=
\left(\abs{z_1}^2-\frac35(\abs{z_1}^2+\abs{z_2}^2)\right)R(z)
\end{equation}
is a homogeneous harmonic quintic.  On $S^3$,
$H_5=H-\frac35R$, and \eqref{eq:QR-annihilation} implies
\[
\Lambda_g(Y_{H_5}(g))=\Lambda_g(Y_H(g)).
\]
The same compact-obstruction argument therefore applies to $Q+\eps H_5$.
Its restriction to $S^3$ is exactly \eqref{eq:main-counterexample}, proving
\cref{thm:main-counterexample}.

\appendix

\section{Proof of algebraic transversality}\label{app:transversality}

We prove \cref{prop:algebraic-transversality}.  Fix $g\in\SO(4)$ for which the
$Q$-labels are exact, write
\[
p_i=gq_i=(a_i,b_i)\in\C^2,
\qquad
a=(a_0,\dots,a_4),\quad b=(b_0,\dots,b_4)\in\C^5,
\]
and use coordinatewise multiplication of vectors in $\C^5$.  The Hermitian
inner product is conjugate-linear in the first variable:
\[
\ip{x}{y}_{\C^5}=\sum_i\overline{x_i}y_i.
\]

\subsection{Frame identities}

The $5\times4$ real matrix whose rows are the $p_i$ has orthonormal columns,
all perpendicular to $\one$.  Therefore
\begin{align}
\sum_i a_i=\sum_i b_i&=0,\label{eq:frame-sums}\\
\sum_i a_i^2=\sum_i b_i^2
=\sum_i a_ib_i=\sum_i a_i\overline b_i&=0,\label{eq:frame-orthog}\\
\sum_i\abs{a_i}^2=\sum_i\abs{b_i}^2&=2,\label{eq:frame-norms}\\
\abs{a_i}^2+\abs{b_i}^2&=\frac45.\label{eq:row-norm}
\end{align}
Consequently
\[
\one,a,b,\overline a,\overline b
\]
are mutually orthogonal in $\C^5$, with squared norms $5,2,2,2,2$.
Every $x\in\C^5$ has the expansion
\begin{equation}\label{eq:basis-expansion}
x=\frac{\sum_i x_i}{5}\one
+\frac12\left(
\ip{a}{x}a+\ip{b}{x}b
+\ip{\overline a}{x}\overline a
+\ip{\overline b}{x}\overline b
\right).
\end{equation}
For $u,v\in\C^5$, write
\[
(u\wedge v)_{ij}=u_iv_j-v_iu_j.
\]

\subsection{Cycle equations for the cubic}

At the root $g\nu_{ij}=(p_i-p_j)/\sqrt2$,
\[
W(g\nu_{ij})
=\frac1{2\sqrt2}(a_i-a_j)^2(\overline b_i-\overline b_j).
\]
Ignoring the harmless factor $1/(2\sqrt2)$, let
\[
\widetilde Y_{ij}=(a_i-a_j)^2(\overline b_i-\overline b_j).
\]
Modulo exact matrices, expansion gives
\begin{equation}\label{eq:cubic-cycle-form}
\Pi\widetilde Y\Pi
\equiv -a^2\wedge\overline b+2a\wedge(a\overline b).
\end{equation}
Indeed, the omitted term is
$a_i^2\overline b_i-a_j^2\overline b_j$.

Define
\begin{align}
s&=\frac12\sum_i\abs{a_i}^2a_i,
&\lambda&=\frac12\sum_i\abs{a_i}^2\overline b_i,\label{eq:s-lambda}\\
m&=\sum_i a_i^2\overline b_i,
&n&=\sum_i a_i^3,
&u&=\frac12\sum_i a_i\overline b_i^2,
&\ell&=\frac12\sum_i a_i^2b_i.\label{eq:m-n-u-ell}
\end{align}
Using \eqref{eq:basis-expansion},
\begin{align}
a^2&=sa+\frac m2b+\frac n2\overline a+\ell\overline b,
\label{eq:a2-expansion}\\
a\overline b&=\lambda a+ub+\frac m2\overline a-s\overline b.
\label{eq:abbar-expansion}
\end{align}
Substitution into \eqref{eq:cubic-cycle-form} yields
\begin{equation}\label{eq:omega}
\omega:=\Pi\widetilde Y\Pi
=2u\,a\wedge b
+m\,a\wedge\overline a
-3s\,a\wedge\overline b
-\frac m2b\wedge\overline b
-\frac n2\overline a\wedge\overline b.
\end{equation}
Exactness of the real $Q$-labels is precisely $\operatorname{Re}\omega=0$.
The six wedges of the complex basis $a,b,\overline a,\overline b$ are linearly
independent.  Comparing \eqref{eq:omega} with its conjugate gives
\begin{equation}\label{eq:cubic-exactness-conditions}
s=0,
\qquad m\in\R,
\qquad 4u=\overline n.
\end{equation}

Replacing $g$ by $h_\theta g$ changes neither the values of $Q$ and $H$ nor
the matrix $X(g)$, while
$\lambda\mapsto e^{-2i\theta}\lambda$.  We may therefore assume
\begin{equation}\label{eq:t-gauge}
\lambda=t\ge0.
\end{equation}
In addition define
\begin{equation}\label{eq:v-k}
v=\sum_i a_ib_i^2,
\qquad
k=\sum_i b_i^3.
\end{equation}

\subsection{The coordinatewise multiplication table}

Expanding each product in the orthogonal basis
$\one,a,b,\overline a,\overline b$ by \eqref{eq:basis-expansion}, and using
\eqref{eq:row-norm} and \eqref{eq:cubic-exactness-conditions}, gives
\begin{align}
a\overline a&=\frac25\one+tb+t\overline b,
&b\overline b&=\frac25\one-tb-t\overline b,\notag\\
a^2&=\frac m2b+\frac n2\overline a+\ell\overline b,
&\overline a^2&=\frac m2\overline b+\frac{\overline n}{2}a+\overline\ell b,
\notag\\
a\overline b&=ta+\frac{\overline n}{4}b+\frac m2\overline a,
&\overline a b&=t\overline a+\frac n4\overline b+\frac m2a,
\notag\\
ab&=ta+\ell\overline a+\frac v2\overline b,
&\overline a\,\overline b&=t\overline a+\overline\ell a
 +\frac{\overline v}{2}b,
\notag\\
b^2&=\frac n4a-tb+\frac v2\overline a+\frac k2\overline b,
&\overline b^2&=\frac{\overline n}{4}\overline a-t\overline b
 +\frac{\overline v}{2}a+\frac{\overline k}{2}b.
\label{eq:multiplication-table}
\end{align}
For example, the coefficient of $b$ in $a\overline a$ is
$\frac12\sum_i\overline b_i\abs{a_i}^2=t$, and the other coefficients are
obtained in the same way.

Coordinatewise multiplication is associative.  Comparing coefficients in
\eqref{eq:multiplication-table} gives the following necessary equations:
\begin{align}
3mn&=4tv,\tag{E1}\label{eq:E1}\\
\overline n\,v+16\ell t+8mt&=0,\tag{E2}\label{eq:E2}\\
-2km+24\ell t-n^2&=0,\tag{E3}\label{eq:E3}\\
-kt+\ell m+4t^2&=0,\tag{E4}\label{eq:E4}\\
\overline n\,k+mn+8tv&=0,\tag{E5}\label{eq:E5}\\
20\abs\ell^2+5m^2+5\abs n^2-40t^2&=8,\tag{E6}\label{eq:E6}\\
16\abs\ell^2-4m^2-\abs n^2+4\abs v^2&=0,\tag{E7}\label{eq:E7}\\
20\abs\ell^2+60t^2+5\abs v^2&=8.\tag{E8}\label{eq:E8}
\end{align}
More precisely, (E1)--(E3) are the $a,b,\overline b$ coefficients of
$(a^2)b=a(ab)$; (E4)--(E5) are the $a,b$ coefficients of
$(ab)b=a(b^2)$; and (E6), (E7), (E8) are respectively the indicated
coefficients of
\[
(a^2)\overline a=a(a\overline a),\qquad
(ab)\overline a=a(b\overline a),\qquad
(ab)\overline b=a(b\overline b).
\]

\subsection{Classification of the possible cubic zeros}

We extract the consequences of (E1)--(E8) needed below.

\medskip
\noindent\textbf{Case I: $t=0$.}
Equation (E1) gives $mn=0$.  If $m=0$, then (E3) gives $n=0$, while
(E6)--(E7) become
\[
20\abs\ell^2=8,
\qquad
16\abs\ell^2+4\abs v^2=0,
\]
a contradiction.  Hence $m\ne0$ and $n=0$.  Equations (E4) and (E3) give
$\ell=k=0$, and then
\begin{equation}\label{eq:case-I}
m^2=\frac85,
\qquad
\abs v^2=m^2.
\end{equation}

\medskip
\noindent\textbf{Case II: $t>0$ and $n=0$.}
Equation (E1) gives $v=0$, and (E2) gives
\begin{equation}\label{eq:case-II-ell}
\ell=-\frac m2.
\end{equation}
Here $m\ne0$, since otherwise \eqref{eq:multiplication-table} would give
$a^2=0$.  Equations (E3) and (E4) give
$k=-6t$ and $m^2=20t^2$.  Finally (E8) yields
\begin{equation}\label{eq:case-II}
t^2=\frac1{20},
\qquad
m^2=1.
\end{equation}

\medskip
\noindent\textbf{Case III: $t>0$ and $n\ne0$.}
First $m\ne0$, since otherwise (E1), (E2), and (E3) successively give
$v=0$, $\ell=0$, and $n=0$.  Put
\begin{equation}\label{eq:q-h}
q=\frac mt,
\qquad
h=\frac{\abs n^2}{t^2}>0.
\end{equation}
Equations (E1) and (E2) give
\begin{equation}\label{eq:v-ell-case-III}
v=\frac{3mn}{4t},
\qquad
\ell=-\frac{m(32+3h)}{64}.
\end{equation}
In particular, $\ell$ is real.  Equation (E5) gives
\[
k=-7m\frac n{\overline n}.
\]
Equation (E4) shows that $k$ is real, so
\begin{equation}\label{eq:sigma}
\sigma:=\frac n{\overline n}\in\{1,-1\},
\qquad
k=-7m\sigma,
\qquad
n^2=\sigma\abs n^2.
\end{equation}
After division by appropriate powers of $t$, equations (E4) and (E3) become
\begin{align}
3hq^2+32q^2-448\sigma q-256&=0,\label{eq:case-III-a}\\
-9hq-8\sigma h+112\sigma q^2-96q&=0.\label{eq:case-III-b}
\end{align}
Eliminating $h$ gives
\begin{align}
(q-4)(q+2)(21q^2+42q+16)&=0 &&(\sigma=1),\label{eq:resultant-plus}\\
(q+4)(q-2)(21q^2-42q+16)&=0 &&(\sigma=-1).
\label{eq:resultant-minus}
\end{align}
The roots $q=-2\sigma$ give $h=-64$ and are impossible.  One root of each
quadratic gives
$h=(-74-14\sqrt{105})/3<0$.  The candidates $q=4\sigma$ give $h=32$;
then \eqref{eq:v-ell-case-III} and (E1) imply
\[
\frac{\abs\ell^2}{t^2}=64,
\qquad
\frac{\abs v^2}{t^2}=288,
\]
and the left side of (E7), divided by $t^2$, is $2080$, a contradiction.
The only remaining possibility is
\begin{equation}\label{eq:case-III-qh}
q=-\sigma\left(1-\frac{\sqrt{105}}{21}\right),
\qquad
h=\frac{-74+14\sqrt{105}}3.
\end{equation}
Equation (E6) then gives
\begin{equation}\label{eq:case-III-tm}
t^2=\frac{7(5\sqrt{105}-33)}{1440},
\qquad
m^2=\frac{49\sqrt{105}-477}{1080}.
\end{equation}
These three cases exhaust every actual cubic zero, because every such zero
satisfies the necessary equations (E1)--(E8).

\subsection{Evaluation of the quintic functional}

For the actual root $g\nu_{ij}=(p_i-p_j)/\sqrt2$,
\begin{equation}\label{eq:quintic-edge}
\abs{z_1}^2W(z)
=\frac1{4\sqrt2}
(a_i-a_j)^3(\overline a_i-\overline a_j)
(\overline b_i-\overline b_j).
\end{equation}
Define
\begin{equation}\label{eq:Sigma}
\Sigma=\sum_{i<j}X(g)_{ij}
(a_i-a_j)^3(\overline a_i-\overline a_j)
(\overline b_i-\overline b_j).
\end{equation}
Then
\begin{equation}\label{eq:Lambda-H-Sigma}
\Lambda_g(Y_H(g))=\frac1{4\sqrt2}\operatorname{Im}\Sigma.
\end{equation}
In wedge notation,
\begin{equation}\label{eq:X-wedge}
X(g)=\frac i2\bigl(\overline a\wedge a+2\overline b\wedge b\bigr).
\end{equation}
Modulo an exact matrix, the degree-five edge expression in
\eqref{eq:Sigma} is
\begin{align}
F\equiv{}&
\overline b\wedge(a^3\overline a)
+\overline a\wedge(a^3\overline b)
-(\overline a\,\overline b)\wedge a^3\notag\\
&+3a\wedge(a^2\overline a\,\overline b)
-3(a\overline b)\wedge(a^2\overline a)\notag\\
&-3(a\overline a)\wedge(a^2\overline b)
+3(a\overline a\,\overline b)\wedge a^2.
\label{eq:F-wedge}
\end{align}
The omitted exact term is
$(a^3\overline a\,\overline b)\wedge\one$, which is annihilated by $X(g)$.
For the bilinear dot product $x\cdot y=\sum_i x_iy_i$,
\begin{equation}\label{eq:wedge-pairing}
\sum_{i<j}(u\wedge v)_{ij}(x\wedge y)_{ij}
=(u\cdot x)(v\cdot y)-(u\cdot y)(v\cdot x).
\end{equation}
Every dot product in \eqref{eq:wedge-pairing} is five times the coefficient of
$\one$ in the corresponding coordinatewise product.  Substituting
\eqref{eq:multiplication-table}, \eqref{eq:X-wedge}, and
\eqref{eq:F-wedge} gives
\begin{equation}\label{eq:Sigma-Braw}
\Sigma=-\frac i{80}\,\mathcal B_{\mathrm{raw}},
\end{equation}
where
\begin{align}
\mathcal B_{\mathrm{raw}}={}&
-960\overline k\,\ell t
+160\abs\ell^2m
+320\ell t^2
-120m^3\notag\\
&-75m\abs n^2
-800mt^2
+100m\abs v^2
+192m
-160nt\overline v.
\label{eq:Braw}
\end{align}
In all three cases above, $\ell$ is real.  The conjugates of (E4) and (E1)
give
\begin{equation}\label{eq:two-substitutions}
\overline k\,t=\ell m+4t^2,
\qquad
nt\overline v=\frac34m\abs n^2.
\end{equation}
Using \eqref{eq:two-substitutions}, followed by (E6) and (E8), reduces
\eqref{eq:Braw} to
\begin{equation}\label{eq:B-simplified}
\mathcal B
=m\bigl(180m^2+105\abs n^2-4400t^2-128\bigr)
-3520\ell t^2.
\end{equation}
The three cases now give
\begin{equation}\label{eq:B-values}
\begin{array}{c@{\qquad}c}
\toprule
\text{case}&\mathcal B\\
\midrule
\mathrm{I}&160m\\
\mathrm{II}&-80m\\
\mathrm{III}&60m(75-7\sqrt{105})\\
\bottomrule
\end{array}
\end{equation}
Each value is nonzero: in every case $m\ne0$, and
$75-7\sqrt{105}>0$.  Since $\mathcal B$ is real,
\eqref{eq:Lambda-H-Sigma} and \eqref{eq:Sigma-Braw} yield
\[
\Lambda_g(Y_H(g))=-\frac{\mathcal B}{320\sqrt2}\ne0.
\]
This proves \cref{prop:algebraic-transversality}.

\begin{remark}[Symbolic verification]
The wedge expansion \eqref{eq:F-wedge}, the extraction of (E1)--(E8), the
resultants \eqref{eq:resultant-plus}--\eqref{eq:resultant-minus}, and the
simplification \eqref{eq:B-simplified} are polynomial identities.  An
accompanying exact-arithmetic SymPy script checks each of them independently.
\end{remark}

\begin{thebibliography}{99}

\bibitem{Chen2026}
Y.~Chen,
\emph{Note on Lebesgue's universal cover problem},
arXiv:2607.27227, 2026.
\url{https://arxiv.org/abs/2607.27227}.

\bibitem{DincaMawhin2021}
G.~Dinc{\u a} and J.~Mawhin,
\emph{Brouwer Degree: The Core of Nonlinear Analysis},
Progress in Nonlinear Differential Equations and Their Applications,
Birkh\"auser, Cham, 2021.
\href{https://doi.org/10.1007/978-3-030-63230-4}{doi:10.1007/978-3-030-63230-4}.

\bibitem{HauselMakaiSzucs2000}
T.~Hausel, E.~Makai, Jr., and A.~Sz\H{u}cs,
\emph{Inscribing cubes and covering by rhombic dodecahedra via equivariant topology},
Mathematika \textbf{47} (2000), 371--397.
\href{https://doi.org/10.1112/S0025579300015965}{doi:10.1112/S0025579300015965}.

\bibitem{JiangLimYaoYe2011}
X.~Jiang, L.-H.~Lim, Y.~Yao, and Y.~Ye,
\emph{Statistical ranking and combinatorial Hodge theory},
Math. Program. \textbf{127} (2011), 203--244.
\href{https://doi.org/10.1007/s10107-010-0419-x}{doi:10.1007/s10107-010-0419-x}.

\bibitem{Karasev2016}
R.~N.~Karasev,
\emph{A note on Makeev's conjectures},
J. Math. Sci. (N.Y.) \textbf{212} (2016), 521--526.
\href{https://doi.org/10.1007/s10958-016-2679-3}{doi:10.1007/s10958-016-2679-3}.

\bibitem{Kuperberg1999}
G.~Kuperberg,
\emph{Circumscribing constant-width bodies with polytopes},
New York J. Math. \textbf{5} (1999), 91--100.
\url{https://nyjm.albany.edu/j/1999/5-6p.pdf}.

\bibitem{Makeev1994}
V.~V.~Makeev,
\emph{Inscribed and circumscribed polyhedra of a convex body},
Math. Notes \textbf{55} (1994), 423--425.
\href{https://doi.org/10.1007/BF02112484}{doi:10.1007/BF02112484}.

\bibitem{MartinVanLoan2008}
C.~D.~Moravitz Martin and C.~F.~Van Loan,
\emph{A Jacobi-type method for computing orthogonal tensor decompositions},
SIAM J. Matrix Anal. Appl. \textbf{30} (2008), 1219--1232.
\href{https://doi.org/10.1137/060655924}{doi:10.1137/060655924}.

\end{thebibliography}

\end{document}
